\chapter{Results and Discussions}

\section{Experimental Setup and Benchmark Datasets}
To evaluate the efficacy, generalizability, and operational viability of the multimodal late-fusion architecture \cite{1}, extensive empirical experiments were conducted across diverse datasets representing both large-scale historical distributions and live zero-day threats.

\subsection{Dataset Profiles}
The experimental evaluation leveraged three primary data sources:
\begin{enumerate}
\item \textbf{Proprietary Enterprise Corpus (3.2 Million URLs):} Compiled from live enterprise proxy logs, security telemetry feeds, and confirmed SOC incident reports over a 12-month period. It comprises $2.8\text{ million}$ verified benign enterprise URLs and $400,000$ verified malicious phishing portals.
\item \textbf{Multimodal Target Lookup Phishing Dataset (MTLP) \cite{5}:} A rigorous, publicly available academic benchmark comprising 60,000 paired multimodal instances (lexical URL, raw HTML DOM tree, and visual page screenshots), balanced across legitimate and phishing samples.
\item \textbf{Live Zero-Day Testbed (100 URLs):} A curated evaluation set of 100 active, unindexed phishing websites collected from real-time honeypots within their initial 30 minutes of deployment, specifically selected to test evasion resilience against commercial feeds.
\end{enumerate}

\subsection{Hardware and Computing Environment}
Model training and distributed inference benchmarking were conducted on an enterprise computing cluster equipped with dual Intel Xeon Platinum processors (64 physical cores, 2.80 GHz), 256 GB of ECC DDR4 RAM, and four NVIDIA A100 Tensor Core GPUs (40 GB VRAM each). The microservice pipeline was deployed within containerized Docker pods managed by a Kubernetes cluster running Ubuntu 22.04 LTS.

\section{Performance Evaluation Metrics}
Model efficacy was evaluated using standard statistical classification metrics:
\begin{itemize}
\item \textbf{Accuracy:} $\text{ACC} = \frac{\text{TP} + \text{TN}}{\text{TP} + \text{TN} + \text{FP} + \text{FN}}$
\item \textbf{Precision:} $\text{PREC} = \frac{\text{TP}}{\text{TP} + \text{FP}}$
\item \textbf{Recall (Sensitivity):} $\text{REC} = \frac{\text{TP}}{\text{TP} + \text{FN}}$
\item \textbf{F1-Score:} $\text{F1} = 2 \times \frac{\text{PREC} \times \text{REC}}{\text{PREC} + \text{REC}}$
\item \textbf{ROC-AUC:} Area under the Receiver Operating Characteristic Curve.
\end{itemize}
Given the severe operational cost of false negatives in cybersecurity, the F1-score serves as the primary benchmark metric.

\section{Multimodal Late Fusion vs Individual Modality Branches}
A critical research inquiry centers on whether combining heterogeneous data modalities provides superior discrimination over individual, specialized models. Table \ref{tab:results_branches} presents the comparative evaluation across each standalone branch and the proposed multimodal late-fusion ensemble.

\begin{table}[htbp]
\centering
\resizebox{\textwidth}{!}{%
\begin{tabular}{|l|c|c|c|c|c|}
\hline
\textbf{Model Architecture / Modality} & \textbf{Accuracy (\%)} & \textbf{Precision} & \textbf{Recall} & \textbf{F1-Score} & \textbf{ROC-AUC} \\
\hline
\multicolumn{6}{|c|}{\textbf{Proprietary Enterprise Corpus (3.2 Million URLs)}} \\
\hline
CatBoost (Tabular URL \& Meta) & 93.8 & 0.942 & 0.928 & 0.935 & 0.971 \\
\hline
1D-CNN (Character-level URL) & 91.2 & 0.915 & 0.898 & 0.906 & 0.954 \\
\hline
CodeBERT (HTML Source Code) & 95.1 & 0.958 & 0.941 & 0.949 & 0.982 \\
\hline
EfficientNet-B7 (Visual Screenshot) & 93.2 & 0.936 & 0.920 & 0.928 & 0.968 \\
\hline
\textbf{Multimodal Late Fusion (Proposed)} & \textbf{99.0} & \textbf{0.991} & \textbf{0.987} & \textbf{0.989} & \textbf{0.998} \\
\hline
\multicolumn{6}{|c|}{\textbf{Public Benchmark Dataset (MTLP) \cite{5}}} \\
\hline
CatBoost (Tabular URL \& Meta) & 89.4 & 0.898 & 0.876 & 0.887 & 0.938 \\
\hline
1D-CNN (Character-level URL) & 86.7 & 0.871 & 0.849 & 0.860 & 0.921 \\
\hline
CodeBERT (HTML Source Code) & 91.5 & 0.921 & 0.904 & 0.912 & 0.956 \\
\hline
EfficientNet-B7 (Visual Screenshot) & 89.8 & 0.905 & 0.887 & 0.896 & 0.944 \\
\hline
\textbf{Multimodal Late Fusion (Proposed)} & \textbf{95.4} & \textbf{0.956} & \textbf{0.950} & \textbf{0.953} & \textbf{0.984} \\
\hline
\end{tabular}%
}
\vspace{0.2cm}
\caption{Comparative Performance of Standalone Branches vs. Multimodal Late Fusion \cite{1}}
\label{tab:results_branches}
\end{table}

As demonstrated in Table \ref{tab:results_branches}, each standalone branch achieves respectable individual performance, with CodeBERT attaining the highest standalone F1-score (0.949 on the enterprise corpus) due to its profound understanding of structural DOM intent. However, the Multimodal Late Fusion meta-classifier dramatically outperforms all single branches, achieving an F1-score of \textbf{0.989} on the enterprise dataset and \textbf{0.953} on the public MTLP benchmark. This verifies that combining orthogonal perspectives (lexical, behavioral, syntactic, visual) successfully neutralizes evasions that bypass single-layer filters.

\section{Zero-Day Phishing Detection Benchmark}
To rigorously evaluate real-world defensive resilience against unindexed threats, the framework was evaluated against 100 live zero-day phishing portals and benchmarked against two widely deployed commercial blacklist feeds: Google Safe Browsing and OpenPhish.

As documented in Table \ref{tab:zeroday}:
\begin{table}[htbp]
\centering
\resizebox{0.85\textwidth}{!}{%
\begin{tabular}{|l|c|c|c|}
\hline
\textbf{Detection Platform} & \textbf{Detected (/100)} & \textbf{Missed (FN)} & \textbf{Evasion Rate (\%)} \\
\hline
Google Safe Browsing (API Lookup) & 41 & 59 & 59.0\% \\
\hline
OpenPhish Automated Feed & 67 & 33 & 33.0\% \\
\hline
\textbf{Proposed Multimodal Late Fusion} & \textbf{92} & \textbf{8} & \textbf{8.0\%} \\
\hline
\end{tabular}%
}
\vspace{0.2cm}
\caption{Zero-Day Detection Benchmark on 100 Live Unindexed Phishing Websites}
\label{tab:zeroday}
\end{table}

\begin{itemize}
\item Google Safe Browsing intercepted only 41 out of 100 zero-day links (a 59\% evasion rate). Because Google Safe Browsing relies primarily on crawler verification and crawler-submitted hash lists, adversaries successfully evaded detection by restricting early traffic through IP geolocation geofencing and dynamic DNS redirection.
\item OpenPhish identified 67 out of 100 links (a 33\% evasion rate), benefiting from automated community submissions but still exhibiting incubation lag.
\item The Proposed Multimodal Architecture correctly identified \textbf{92 out of 100} live attacks (an evasion rate of only 8.0\%), successfully intercepting zero-day links minutes after registration by identifying underlying structural DOM anomalies and visual brand mimicry.
\end{itemize}

\section{SOC Analyst Usability and Interpretability Evaluation}
A distinctive contribution of this research is evaluating whether Explainable AI actually improves real-world human analyst efficiency. A user study was conducted involving 25 tier-1 and tier-2 SOC security analysts tasked with triaging 50 ambiguous alerts across three explanation presentation modes:
\begin{enumerate}
\item \textbf{Raw Quantitative XAI:} Feature importance bar charts, Grad-CAM heatmaps, and raw token attention scores.
\item \textbf{Local LLM Synthesized Report:} Natural language incident summary generated from XAI features.
\item \textbf{RAG-Augmented Local LLM Report:} Natural language summary enriched with MITRE ATT\&CK techniques and historical campaign attribution.
\end{enumerate}

Analysts scored each mode across four criteria on a 5-point Likert scale:
\begin{enumerate}
\item \textbf{Comprehensibility:} Ease of understanding the rationale without data science expertise.
\item \textbf{Actionability:} Clarity regarding whether to block, isolate, or whitelist the domain.
\item \textbf{Trust and Confidence:} Degree of confidence in the model's recommendation.
\item \textbf{Diagnostic Speed:} Perceived reduction in verification and triage effort.
\end{enumerate}

As shown in Table \ref{tab:usability}:
\begin{table}[htbp]
\centering
\resizebox{\textwidth}{!}{%
\begin{tabular}{|p{4.8cm}|c|c|c|c|c|}
\hline
\textbf{Explanation Modality} & \textbf{Comprehens.} & \textbf{Actionab.} & \textbf{Trust} & \textbf{Speed} & \textbf{Overall Mean} \\
\hline
Raw XAI Plots (SHAP, Grad-CAM, IG) & 3.12 & 3.35 & 3.60 & 3.61 & 3.42 \\
\hline
Local LLM Natural Language Report & 4.02 & 3.94 & 3.88 & 4.04 & 3.97 \\
\hline
\textbf{RAG-Augmented Local LLM Report} & \textbf{4.28} & \textbf{4.18} & \textbf{4.15} & \textbf{4.19} & \textbf{4.20} \\
\hline
\end{tabular}%
}
\vspace{0.2cm}
\caption{SOC Analyst Usability and Interpretability Evaluation Scores (Scale: 1--5)}
\label{tab:usability}
\end{table}

\begin{itemize}
\item Presenting Raw XAI outputs (bar charts and heatmaps) scored lowest ($3.42/5$), as analysts reported cognitive fatigue in cross-referencing multiple disparate graphs.
\item Synthesizing explanations into natural language briefs elevated usability to $3.97/5$.
\item Adding RAG contextualization achieved the highest rating (\textbf{4.20/5}), reducing mean analyst verification time from $4.5\text{ minutes}$ down to under $45\text{ seconds}$ per incident.
\end{itemize}

\section{System Latency and Microservice Throughput}
In production enterprise environments, defensive models must not introduce unacceptable browsing delays. As detailed in Table \ref{tab:latency}, because the data collection and model inference branches execute concurrently within an asynchronous containerized mesh, the overall pipeline latency is determined by the slowest branch (headless browser screenshot rendering, taking $84.6\text{ ms}$).

\begin{table}[htbp]
\centering
\resizebox{\textwidth}{!}{%
\begin{tabular}{|l|c|l|}
\hline
\textbf{Pipeline Component} & \textbf{Mean Latency (ms)} & \textbf{Execution Paradigm} \\
\hline
URL Feature \& Metadata Extraction & 12.4 & Concurrent Async Worker \\
\hline
1D-CNN Character-level Inference & 7.8 & GPU Parallelized \\
\hline
Headless Browser Crawl \& Screenshot & 84.6 & Asynchronous Chrome Pod \\
\hline
CodeBERT DOM Tokenization \& Inference & 64.2 & GPU Tensor Batch \\
\hline
EfficientNet-B7 Visual Inference & 38.5 & GPU Tensor Batch \\
\hline
Late Fusion Meta-Classifier & 1.8 & CPU Lightweight Meta \\
\hline
Local LLM Incident Synthesis (Quantized) & 185.0 & Triggered on High Severity \\
\hline
\textbf{Total End-to-End Classification} & \textbf{108.6} & \textbf{Sub-second Streaming Pipeline} \\
\hline
\end{tabular}%
}
\vspace{0.2cm}
\caption{Component-Wise Execution Latency and Computational Overhead}
\label{tab:latency}
\end{table}

The total classification decision is reached in an average of \textbf{108.6 ms}, fully meeting real-world enterprise proxy inline inspection requirements. The LLM synthesis branch ($185\text{ ms}$) is invoked asynchronously only when the meta-classifier flags a high-severity threat, avoiding unnecessary compute overhead for benign traffic.
